\documentclass[12pt]{article}

\usepackage[thinlines]{easytable}

\usepackage{amsmath}
\usepackage{amssymb}
\usepackage{amsthm}
\usepackage{fullpage}

\newcommand{\R}{\mathbb{R}}
\newcommand{\Requests}{\mathfrak{R}}

\DeclareMathOperator*{\argmax}{argmax}
\DeclareMathOperator*{\argmin}{argmin}
\DeclareMathOperator*{\Exp}{\mathbf{E}}

\DeclareMathOperator*{\maximize}{maximize}
\DeclareMathOperator*{\minimize}{minimize}
\DeclareMathOperator*{\subjectto}{subject\ to}

\newtheorem{theorem}{Theorem}
\newtheorem{lemma}[theorem]{Lemma}
\newtheorem{proposition}[theorem]{Proposition}
\newtheorem{conjecture}[theorem]{Conjecture}

\theoremstyle{definition}
\newtheorem{definition}[theorem]{Definition}

\begin{document}

\title{Pacing}
\author{D\'avid P\'al}
\date{June 20, 2026}

\maketitle

\begin{abstract}
We show convergence results for the pacing problem.
\end{abstract}

\section{Introduction}

Consider a sequence of requests $R_1, R_2, \dots, R_T \in \Requests$. For each
request we have to make one of $N$ decisions. We make a simple assumption that
the requests are sampled i.i.d. from a fixed, but unknown, distribution $D$.
Each request is an $N \times (M+1)$ matrix with entries in the interval
$[-1,1]$.\footnote{Any bounded interval would work. We chose $-1$ and $1$ for
convenience.} In other words, we assume that the space of requests is
$\Requests = [-1,1]^{N \times (M+1)}$. We would like to come up with a
(randomized) policy $\pi$ that maximizes a linear reward subject to linear
constraints. Formally, a randomized policy $\pi$ is a function $\pi:\Requests
\to \Delta_N$, where
$$
\Delta_N = \left\{ x \in \R^N ~:~ x_1, x_2, \dots, x_N \ge 0, \sum_{i=1}^N x_i = 1  \right\}
$$
is the probability simplex. We would like to come up with a policy that is an
``close'' to an optimal solution of the optimization problem
\begin{align*}
\maximize_{\pi \in \Pi} \qquad & \frac{1}{T} \sum_{t=1}^T \sum_{i=1}^N R_{t,i,0} \cdot \pi(R_t)_i \\
\subjectto \qquad & \frac{1}{T} \sum_{t=1}^T \sum_{i=1}^N R_{t,i,j} \cdot \pi(R_t)_i \le B_j & \text{for all $j=1,2,\dots,M$}
\end{align*}
The maximization is over the set $\Pi$ of all set of measurable functions
$\pi:\Requests \to \Delta_N$ with respect to the $\sigma$-algebras defined by
the Borel sets of $\Requests$ and $\Delta_N$ respectively.

The value $R_{t,i,0}$ is the reward collected for $t$-th request if we make
decision $i$ for it. The value $R_{t,i,j}$ for $j = 1, 2, \dots, M$ is the
$j$-th business metric for the $t$-th request if we make decision $i$ for it.
The right-hand sides of the constraints $B_1, B_2, \dots, B_M$ are given to us
ahead of time.

For example, at Uber, a request can be a request by a customer to get a ride and
Uber needs to choose one of $K$ possible prices for the request. $R_{t,i,0}$
could be the probability of the trip happening, $R_{t,i,1}$ could be the price
paid by the rider, $R_{t,i,2}$ could be the price paid to the driver, etc. The
right-hand sides of the constraints are the business objectives Uber wishes to
achieve over the period of say, one day or one week.

The optimization problem can be expressed as a standard linear program:

\begin{align*}
\maximize_{x_{1,1}, \dots, x_{T,N} \in \R} \qquad & \frac{1}{T} \sum_{t=1}^T \sum_{i=1}^N R_{t,i,0} \cdot x_{t,i} \\
\subjectto \qquad & \frac{1}{T} \sum_{t=1}^T \sum_{i=1}^N R_{t,i,j} \cdot x_{t,i} \le B_j & \text{for all $j=1,2,\dots,M$} \\
& \sum_{i=1}^N x_{t,i} = 1 & \text{for all $t=1,2,\dots,T$} \\
& x_{t,i} \ge 0 & \text{for all $i=1,2,\dots,N$ and all $t=1,2,\dots,T$}
\end{align*}

However, as opposed to a standard linear program, we don't have all the
coefficients $R_{t,i,j}$ available upfront. Instead, the requests come one at a
time, and for each request we need to choose the probability vector $(x_{t,1},
x_{t,2}, \dots, x_{t,N}) \in \Delta_N$. In order to make the decision for $t$-th
request quickly, the probability vector is computed on the fly using the policy
$\pi$ applied to the request $R_t$.

We focus on a regime where the number of requests $T$ is much greater than both
$N$ and $M$. The assumption that the entries of a request matrix $R_t$ lie in
$[-1,1]$ translates to an assumption that each request has negligible impact on
the total (or average) business metrics. For the applications that we are
interested in, the i.i.d. assumption is a good approximation for short enough
time periods (e.g. a day or a week). Perhaps the most problematic assumption is
that the value of the reward and business metrics $R_{t,i,j}$ are known at the
time when the decision $\pi(R_t)$ needs to be made. In applications we have in
mind, the reward and the business metrics are predicted by machine learning
models at time of the decision, and the actual reward and other business metrics
are revealed only after the decision has been made, sometimes with a significant
delay. However, in order to keep our model simple, we assume that $R_t$ is
available at the time when the decision for $t$-th request needs to be made, and
we assume that the entries in $R_t$ are perfectly accurate.

We also make the simplifying assumption that the number of requests $T$ is known
upfront. In applications, this value would be predicted by a machine learning
model.

We use the following notation for the rest of the paper. For $i=1,2,\dots,N$, we
define $e_i$ to be the vector in $\R^N$ with $i$-th entry equal to $1$ and all
other entries equal to zero, that is,
$$
e_i = \underbrace{(0,\dots,0,1,0,\dots,0)}_{\text{$1$ is at position $i$}} \: .
$$

\section{Lagrangian policies}

We consider policies parametrized by $M+1$ real numbers $\lambda_0, \lambda_1,
\dots, \lambda_M$. As will become clear later, the $\lambda_1, \dots, \lambda_M$
can be interpreted as the dual variables associated with the constraints of the
linear program. However, for now, they are arbitrary real numbers. For any
$(M+1)$-dimnesional vector $\lambda = (\lambda_0, \lambda_1, \dots, \lambda_M)
\in \R^{M+1}$ we define a policy $\pi_\lambda:\Requests \to \Delta_N$ as
follows. For any request $R \in \Requests$, we define
$$
\pi_\lambda(R) = e_{i^*}
\qquad \text{where} \qquad
i^* = \argmax_{i=1,2,\dots,N} \lambda_0 R_{i,0} - \lambda_1 R_{i,1} - \lambda_2 R_{i,2} - \dots - \lambda_M R_{i,M} \: .
$$
The ties in the $\argmax$ are broken towards the smaller index $i$. Obviously,
the policy $\pi_\lambda$ is deterministic and measurable. We call a policy of
this form a \emph{deterministic Lagrangian policy}, as it corresponds to a
Lagrangian relaxation of the linear program.

We denote by $\Pi_\Lambda^\mathrm{det}$ the set of deterministic Lagrangian
policies corresponding to nonzero vectors $\lambda \in \R^{M+1}$. Formally, that is,
$$
\Pi_\Lambda^\mathrm{det} = \left\{\pi_\lambda ~:~ \lambda \in \R^{M+1}, \lambda \neq 0 \right\} \: .
$$
We exclude the all-zero vector $\lambda$ because it is the point of
discontinuity for a certain function that we will define later.

For any positive number $c > 0$ and vector $\lambda \in \R^{M+1}$, the policies
$\pi_{c \lambda}$ and $\pi_\lambda$ are the same. We can thus assume without
loss of generality that $\|\lambda\|_2 = 1$. Thus,
$$
\Pi_\Lambda^\mathrm{det} = \left\{\pi_\lambda ~:~ \lambda \in \R^{M+1}, \|\lambda\|_2 = 1 \right\} \: .
$$
Later we will use that fact that the set $\{ \lambda \in \R^{M+1} ~:~  \|\lambda\|_2 = 1\}$
is compact.

\section{Failure of deterministic policies}

In general, determistic policies are not good solutions for certain sequences of
requests. In this section, we provide the simplest possible example that
illustrates the point. Suppose $N=2$ and $M=1$ and $T$ can be atrbitrary. For
all requests,
$$
R_t
= \begin{pmatrix} R_{t,1,0} & R_{t,1,1} \\ R_{t,2,0} & R_{t,2,1} \end{pmatrix}
= \begin{pmatrix} 1 & 1 \\ 0 & 0 \end{pmatrix}
$$
and $B_1 = \frac{1}{2}$. Any optimal policy $\pi^*$ for this sequence of
requests satisfies
\begin{equation}
\label{equation:uniform-random-policy-on-two-outcomes}
\pi^*(R_t) = \left( \frac{1}{2}, \frac{1}{2} \right) \: .
\end{equation}
The output of a policy for requests other than $R_t$ does not play any role and
it can be arbitrary. More importantly, any policy that does not satisfy
condition \eqref{equation:uniform-random-policy-on-two-outcomes} is suboptimal.

However, as we will see later, deterministic Lagrangian policies are almost
optimal assuming that the requests are sampled i.i.d. from a distribution with a
probability density.

\section{Uniform convergence for Lagrangian policies}

\begin{theorem}
Let $D$ be any distribution over $\Requests$. Let $\delta \in (0,1)$ and $T$ be
a positive integer. Suppose $R, R_1, R_2, \dots, R_T$ is an i.i.d. sample from
$D$. Then, with probability at least $1-\delta$, for all $\lambda \in \R^{M+1}$ and for all $j=0,1,2,\dots,M$,
$$
\left| \Exp_{R \sim D} \left[ \sum_{i=1}^N \pi_\lambda(R)_i \cdot R_{i,j}  \right]  -  \frac{1}{T} \sum_{t=1}^T \sum_{i=1}^N R_{t,i,j} \cdot \pi_\lambda(R_t)_i  \right|
\le \sqrt{\frac{8(M+1) (\ln(4T) + 2 \ln N)}{T}} + \sqrt{ \frac{2\log(1/\delta)}{T} } \: .
$$
\end{theorem}

\begin{proof}
Let $R, R_1, R_2, \dots, R_T$ be an i.i.d. sample from $D$ of size $T+1$.
Consider the random variable
$$
Z = \sup_{\substack{\lambda \in \R^{M+1} \\ j=0,1,2,\dots,M}}
\left| \Exp_{R \sim D} \left[ \sum_{i=1}^N \pi_\lambda(R)_i \cdot R_{i,j}  \right] - \frac{1}{T} \sum_{t=1}^T \sum_{i=1}^N R_{t,i,j} \cdot \pi_\lambda(R_t)_i \right| \: .
$$
The random variable $Z$ is a function of $R_1, R_2, \dots, R_T$. It satisfies
the bounded difference property. Namely, if a single $R_t$ changes, then $Z$
changes by at most $2/T$. McDiarmid's inequality implies that for any $\epsilon > 0$,
$$
\Pr\left[ Z \ge \Exp[Z] + \epsilon \right] \le e^{-T\epsilon^2/2} \: .
$$
We set $\delta = e^{-T\epsilon^2/2}$ and solve for $\epsilon$. We get that with
probability at least $1-\delta$,
$$
Z \le \Exp[Z] + \sqrt{ \frac{2\log(1/\delta)}{T} } \: .
$$
It remains to upper bound $\Exp[Z]$. We use the symmetrization trick. Let $R'_1,
R'_2, \dots, R'_T$ be an i.i.d. sample from $D$ of size $T$. Let $\sigma_1,
\sigma_2, \dots, \sigma_T$ be i.i.d. Rademacher random variables, i.e.,
$\Pr[\sigma_t = 1] = \Pr[\sigma_t = -1] = 1/2$ for all $t=1,2,\dots,T$. We
assume that $\sigma_1, \sigma_2, \dots, \sigma_T$, $R$, $R_1, R_2, \dots, R_T$,
$R'_1, R'_2, \dots, R'_T$ are mutually independent. We have
\begin{align*}
\Exp[Z] & = \Exp\left[\sup_{\substack{\lambda \in \R^{M+1} \\ j=0,1,2,\dots,M}} \left| \Exp \left[ \sum_{i=1}^N \pi_\lambda(R)_i \cdot R_{i,j}  \right] - \frac{1}{T} \sum_{t=1}^T \sum_{i=1}^N R_{t,i,j} \cdot \pi_\lambda(R_t)_i \right| \right] \\
& = \Exp\left[\sup_{\substack{\lambda \in \R^{M+1} \\ j=0,1,2,\dots,M}} \left| \Exp \left[ \frac{1}{T} \sum_{t=1}^T \sum_{i=1}^N \pi_\lambda(R'_t)_i \cdot R'_{t,i,j}  \right] - \frac{1}{T} \sum_{t=1}^T \sum_{i=1}^N R_{t,i,j} \cdot \pi_\lambda(R_t)_i \right| \right] \\
& \le \Exp\left[\sup_{\substack{\lambda \in \R^{M+1} \\ j=0,1,2,\dots,M}} \left| \frac{1}{T} \sum_{t=1}^T \sum_{i=1}^N \pi_\lambda(R'_t)_i \cdot R'_{t,i,j} - \frac{1}{T} \sum_{t=1}^T \sum_{i=1}^N R_{t,i,j} \cdot \pi_\lambda(R_t)_i \right| \right] \\
& = \frac{1}{T} \Exp\left[\sup_{\substack{\lambda \in \R^{M+1} \\ j=0,1,2,\dots,M}} \left| \sum_{t=1}^T \sum_{i=1}^N \sigma_t \left(\pi_\lambda(R'_t)_i \cdot R'_{t,i,j} - R_{t,i,j} \cdot \pi_\lambda(R_t)_i \right) \right| \right] \\
& \le \frac{1}{T} \Exp\left[\sup_{\substack{\lambda \in \R^{M+1} \\ j=0,1,2,\dots,M}} \left| \sum_{t=1}^T \sum_{i=1}^N \sigma_t \cdot \pi_\lambda(R'_t)_i \cdot R'_{t,i,j} \right| + \left| \sum_{t=1}^T \sum_{i=1}^N \sigma_t \cdot \pi_\lambda(R_t)_i \cdot R_{t,i,j} \right|  \right] \\
& \le \frac{2}{T} \Exp\left[\sup_{\substack{\lambda \in \R^{M+1} \\ j=0,1,2,\dots,M}} \left| \sum_{t=1}^T \sum_{i=1}^N \sigma_t \cdot \pi_\lambda(R_t)_i \cdot R_{t,i,j}  \right| \right]
\end{align*}

Condition on the event $R_1=r_1, R_2=r_2, \dots, R_T=r_T$ and consider the set of $T$-dimensional vectors
$$
V = \left\{ \left(\sum_{i=1}^N \pi_\lambda(r_1)_i \cdot r_{1,i,j}, \dots, \sum_{i=1}^N \pi_\lambda(r_T)_i \cdot r_{T,i,j} \right) ~:~ \lambda \in \R^{M+1}, j=0,1,2,\dots,M \right\} \: .
$$
The set $V$ is finite. We reduce the supremum over an uncountably infinite set to a maximum over the finite set $V$,
$$
\Exp\left[\sup_{\substack{\lambda \in \R^{M+1} \\ j=0,1,2,\dots,M}} \left| \sum_{t=1}^T \sum_{i=1}^N \sigma_t \cdot \pi_\lambda(R_t)_i \cdot R_{t,i,j}  \right| ~\middle|~ R_1=r_1, R_2=r_2, \dots, R_T=r_T \right]
= \Exp\left[\max_{v \in V} \left| \sum_{t=1}^T \sigma_t v_t \right| \right]  \: .
$$
Consider any vector $v = (v_1, v_2, \dots, v_T)$ in $V$. Any entry $v_t$ of $v$ satisfies
$$
|v_t| = \left|\sum_{i=1}^N \pi(r_t)_i \cdot r_{t,i,j}\right| \le \left(\sum_{i=1}^N \pi(r_t)_i\right) \cdot  \max_{i=1,2,\dots,N} |r_{t,i,j}| \le 1 \: .
$$
This implies that the random variable $Y_v = \sum_{t=1}^T \sigma_t v_t$ is a
$\sigma$-sub-Gaussian random variable with $\sigma=\sqrt{T}$. That is, $Y_v$
satisfies
$$
\Exp[e^{sY_v}] \le e^{s^2T/2} \qquad \text{for all $s \in \R$.}
$$
Let $K$ be such that $|V| \le K$ for any choice $r_1, r_2, \dots, r_T \in
\Requests$. Lemma~\ref{lemma:maximum-of-sub-gaussian-variables}
implies that
$$
\Exp\left[\max_{v \in V} \left| \sum_{t=1}^T \sigma_t v_t \right| \right] \le \sqrt{2T \log(2K)}
$$
Thus, taking expectations over $R_1, R_2, \dots, R_T$, we get that
$$
\Exp\left[\sup_{\substack{\lambda \in \R^{M+1} \\ j=0,1,2,\dots,M}}  \left| \sum_{t=1}^T \sum_{i=1}^N \sigma_t \cdot \pi_\lambda(R_t)_i \cdot R_{t,i,j}  \right|  \right] \le \sqrt{2T \log(2K)} \: .
$$
Therefore,
$$
\Exp[Z] \le \sqrt{\frac{8 \log(2K)}{T}} \: .
$$

It remains to find an upper bound $K$ on the cardinality of the set $V$. Let
$$
U = \left\{ \left(\pi_\lambda(r_1), \dots, \pi_\lambda(r_T) \right) ~:~ \lambda \in \R^{M+1} \right\} \: .
$$
Clearly, $|V| \le (M+1) |U|$. It remains to upper bound the cardinality of $U$.

Consider any $\lambda \in \R^{M+1}$ and the corresponding deterministic
Lagrangian policy $\pi_\lambda$. The output of the policy $\pi_\lambda$ on a
request $r_t$ is determined by $\binom{N}{2}$ indicators
$$
\mathbf{1} \left[\lambda_0 r_{t,i,0} - \lambda_1 r_{t,i,1} - \dots - \lambda_M r_{t,i,M} \ge \lambda_0 r_{t,i',0} - \lambda_1 r_{t,i',1} - \dots - \lambda_M r_{t,i',M}\right]
$$
indexed by pairs $i,i'$ such that $1 \le i < i' \le N$. Using vector notation, we can write the indicator as
$$
\mathbf{1} \left[ \lambda^T (r_{t,i} - r_{t,i'}) \ge 0 \right] \: .
$$
Let
$$
\kappa_\lambda(r_t) = (\mathbf{1} \left[ \lambda^T (r_{t,1} - r_{t,2}) \ge 0 \right], \mathbf{1} \left[ \lambda^T (r_{t,1} - r_{t,3}) \ge 0 \right], \dots, \mathbf{1} \left[ \lambda^T (r_{t,N-1} - r_{t,N}) \ge 0 \right])
$$
be the vector of these $\binom{N}{2}$ indicators. Let
$$
W = \left\{ (\kappa_\lambda(r_1), \kappa_\lambda(r_2), \dots, \kappa_\lambda(r_T)) ~:~ \lambda \in \R^{M+1} \right\} \: .
$$
Clearly, $|U| \le |W|$. It remains to upper bound the cardinality of $W$. The
cardinality of $W$ is upper bounded by the number of behaviors of homogeneous
hyperplanes in $\R^{M+1}$ on $\binom{N}{2} T$ points. The Vapnik-Chervonenkis
dimension of homogeneous hyperplanes in $\R^{M+1}$ is $M+1$.\footnote{The
Vapnik-Chervonenkis dimension is not affacted by the inclusion of the constant $1$ predictor that corresponds to $\lambda = 0$.}
Thus, the number of behaviors of homogeneous hyperplanes in $\R^{M+1}$ on
$\binom{N}{2} T$ points and thus the cardinality of $W$ is at most
$$
|W| \le \sum_{k=0}^{M+1} \binom{T \binom{N}{2}}{k} \le \left(T \binom{N}{2} + 1 \right)^{M+1} \le (N^2T)^{M+1} \: .
$$
Therefore, the cardinality of $V$ is upper bounded as
$$
|V| \le (M+1) |W| \le (M+1) \left(N^2T \right)^{M+1} \le 2^{M+1} \left(N^2T \right)^{M+1} = \left(2N^2T \right)^{M+1} \: .
$$
Using the value $K=\left(2N^2T \right)^{M+1}$, we get that
$$
\Exp[Z] \le \sqrt{\frac{8 \ln(2K)}{T}} \le \sqrt{\frac{8 \ln((4N^2T)^{M+1})}{T}} = \sqrt{\frac{8(M+1) (\ln(4T) + 2 \ln N)}{T}} \: .
$$
\end{proof}


\section{Linear programs}

In order to show that Lagrangian policies are close to the policy, we study the
linear program
\begin{align*}
\maximize_{x_{1,1}, \dots, x_{T,N} \in \R} \qquad & \frac{1}{T} \sum_{t=1}^T \sum_{i=1}^N R_{t,i,0} \cdot x_{t,i} \\
\subjectto \qquad & \frac{1}{T} \sum_{t=1}^T \sum_{i=1}^N R_{t,i,j} \cdot x_{t,i} \le B_j & \text{for all $j=1,2,\dots,M$} \\
& \sum_{i=1}^N x_{t,i} = 1 & \text{for all $t=1,2,\dots,T$} \\
& x_{t,i} \ge 0 & \text{for all $i=1,2,\dots,N$ and all $t=1,2,\dots,T$}
\end{align*}
Its dual is
\begin{align*}
\minimize_{\substack{\mu_1, \dots, \mu_T \in \R \\ \lambda_1, \dots, \lambda_M \in \R}} \qquad & \sum_{t=1}^T \mu_t + \sum_{j=1}^M B_j \lambda_j \\
\subjectto \qquad
& \mu_t + \frac{1}{T} \sum_{j=1}^M R_{t,i,j} \lambda_j  \ge \frac{1}{T} R_{t,i,0} & \text{for all $i=1,2,\dots,N$ and all $t=1,2,\dots,T$} \\
& \lambda_j \ge 0 & \text{for all $j=1,2,\dots,M$}
\end{align*}
We prove several results about these linear programs.

\begin{proposition}[Feasibility]
\label{proposition:feasibility}
The primal program is bounded. If the primal program is feasible, the dual
program is feasible and both programs have an optimal solution.
\end{proposition}

\begin{proposition}[Computing $\mu$]
\label{proposition:computing_mu}
Let $\mu_1, \dots, \mu_T$, $\lambda_1, \dots, \lambda_M$ be an optimal solution of the dual program.
Then,
$$
\mu_t = \frac{1}{T} \max_{i=1,2,\dots,N} \left( R_{t,i,0} - \sum_{j=1}^M R_{t,i,j} \lambda_j \right) \: .
$$
\end{proposition}

\begin{proof}
The dual constraints imply that
$$
\mu_t \ge \frac{1}{T} \left( R_{t,i,0} - \sum_{j=1}^M R_{t,i,j} \lambda_j \right) \: .
$$
Taking maximum over $i$, we get
$$
\mu_t \ge \frac{1}{T} \max_{i=1,2,\dots,N} \left( R_{t,i,0} - \sum_{j=1}^M R_{t,i,j} \lambda_j \right) \: .
$$
Any value of $\mu_t$ that satisfies this inequality is a feasible solution of
the dual program. Since $\mu_t$ appears with coefficients $1$ in the dual
objective function, the inequality must hold with equality.
\end{proof}

\begin{proposition}[Computing primal from dual]
\label{proposition:computing_primal_from_dual}
Let $\mu_1, \dots, \mu_T$, $\lambda_1, \dots, \lambda_M$ be an optimal solution of the dual program.
Let $x_{1,1}, \dots, x_{t,N}$ be the optimal solution of the primal program.
For any index $t$, if
$$
i(t) = \argmax_{i=1,2,\dots,N} \left( R_{t,i,0} - \sum_{j=1}^M R_{t,i,j} \lambda_j \right)
$$
is uniquely defined, then $x_{t,i(t)} = 1$ and $x_{t,i} = 0$ for all $i \neq i(t)$.
\end{proposition}

\begin{proof}
Suppose
$$
i(t) = \argmax_{i=1,2,\dots,N} \left( R_{t,i,0} - \sum_{j=1}^M R_{t,i,j} \lambda_j \right)
$$
is uniquely defined. Proposition~\ref{proposition:computing_mu} implies that
$$
\mu_t = \frac{1}{T} \max_{i=1,2,\dots,N} \left( R_{t,i,0} - \sum_{j=1}^M R_{t,i,j} \lambda_j \right) = \frac{1}{T} \left( R_{t,i(t),0} - \sum_{j=1}^M R_{t,i(t),j} \lambda_j \right) \: .
$$
These two facts together imply that
$$
\mu_t > \frac{1}{T} \left( R_{t,i,0} - \sum_{j=1}^M R_{t,i,j} \lambda_j \right) \qquad \text{for all $i \neq i(t)$.}
$$
Equivalently,
$$
\mu_t + \frac{1}{T} \sum_{j=1}^M R_{t,i,j} \lambda_j > \frac{1}{T} R_{t,i,0} \qquad \text{for all $i \neq i(t)$.}
$$
By complementary slackness,
$$
x_{t,i} = 0 \qquad \text{for all $i \neq i(t)$}.
$$
The primal constraint $\sum_{i=1}^N x_{t,i} = 1$ then implies that $x_{t,i(t)} = 1$.
\end{proof}

Proposition~\ref{proposition:computing_mu} allows us to compute $\mu_1, \dots,
\mu_T$ from optimal $\lambda_1, \dots, \lambda_M$.
Proposition~\ref{proposition:computing_primal_from_dual} allows us to compute
some values of $x_{t,i}$.

\begin{lemma}[Ties are rare]
Suppose that the requests $R_1, R_2, \dots, R_T$ are sampled i.i.d. from a
probability distribution $D$ with a probability density (with respect to the
Lebesgue measure on $\Requests$). Let $\mu_1, \dots, \mu_T$, $\lambda_1, \dots,
\lambda_M$ be an optimal solution of the dual program. With probability one, for
at least $T-M$ values of the index $t$, the argmax
$$
\argmax_{i=1,2,\dots,N} \left( R_{t,i,0} - \sum_{j=1}^M R_{t,i,j} \lambda_j \right)
$$
is uniquely defined.
\end{lemma}

\begin{proof}
TODO
\end{proof}

\begin{lemma}[Lagrangian policies are close to optimal]
Let $D$ be a probability density (with respect to Lebesgue measure on
$\Requests$). Let $R_1, R_2, \dots, R_T$ be sampled i.i.d. from $D$. Let
$x_{1,1}, \dots, x_{T,N}$ be any optimal solution of the primal program. Let
$\delta \in (0,1)$. With probability at least $1-\delta$, there exists $\lambda
\in \R^{M+1}$ such that $\|\lambda\|_2 = 1$ and for all $j=0,1,\dots,M$,
$$
\frac{1}{T} \left| \sum_{t=1}^T \sum_{i=1}^N R_{t,i,j} \left( \pi_\lambda(R_t) - x_{t,i} \right) \right|
\le
\frac{M}{T} + \sqrt{\frac{8(M+1) (\ln(4T) + 2 \ln N)}{T}} + \sqrt{ \frac{2\log(1/\delta)}{T} } \: .
$$
\end{lemma}

\begin{proof}
TODO
\end{proof}

\section{Optimization problem for the expected values}

Recall that $D$ is a distribution over requests. We assume that a request $R$ is
sampled from $D$ at random. We consider the following optimization problem over all possible
measurable policies $\pi:\Requests \to \Delta_N$.

\begin{align*}
\maximize_{\pi} \qquad & \Exp_{R \sim D} \left[ \sum_{i=1}^N \pi(R)_i \cdot R_{i,0} \right] \\
\subjectto \qquad & \Exp_{R \sim D} \left[ \sum_{i=1}^N \pi(R)_i \cdot R_{i,j} \right] \le B_j & \text{for all $j=1,2,\dots,M$} \\
\end{align*}

\begin{theorem}[Lagrangian policies are optimal]
Let $D$ be a probability density (with respect to Lebesgue measure on $\Requests$).
Then, there exists $\lambda \in \R^{M+1}$ such that $\|\lambda\|_2 = 1$
and $\pi_\lambda$ is an optimal solution of the optimization problem.
\end{theorem}

\begin{proof}
TODO
\end{proof}

\appendix

\section{Technical lemmas}

\begin{lemma}[Maximum of sub-Gaussian variables]
\label{lemma:maximum-of-sub-gaussian-variables}
Let $\sigma \ge 0$ and let $K \ge 1$. Let $Y_1, Y_2, \dots, Y_K$ be random
variables that satisfy the $\sigma$-sub-Gaussian property
$$
\Exp\left[ e^{sY_k} \right] \le e^{s^2\sigma^2/2} \qquad \text{for all $s \in \R$ and for all $k=1,2,\dots,K$.}
$$
Then,
$$
\Exp\left[ \max_{k=1,2,\dots,K} |Y_k| \right] \le \sigma \sqrt{2 \ln(2K)} \: .
$$
\end{lemma}

\begin{proof}
For any $s \in \R$,
\begin{align*}
\exp\left(s \Exp\left[ \max_{k=1,2,\dots,K} |Y_k| \right] \right)
& = \exp\left(s \Exp\left[ \max_{k=1,2,\dots,K} \max \left\{Y_k, -Y_k\right\} \right] \right) \\
& \le \Exp\left[ \exp\left(s \max_{k=1,2,\dots,K} \max \left\{Y_k, -Y_k\right\} \right) \right] \\
& = \Exp\left[ \max_{k=1,2,\dots,K} \max \left\{e^{sY_k}, e^{-sY_k}\right\} \right] \\
& \le \Exp\left[ \sum_{k=1}^K \left( e^{sY_k} + e^{-sY_k} \right) \right] \\
& = \sum_{k=1}^K \left( \Exp\left[ e^{sY_k} \right] + \Exp\left[ e^{-sY_k} \right] \right) \\
& \le \sum_{k=1}^K \left( e^{s^2\sigma^2/2} + e^{s^2\sigma^2/2} \right) \\
& = 2K e^{s^2\sigma^2/2} \: .
\end{align*}
The first equality follows from that $|Y_k| = \max \left\{Y_k, -Y_k\right\}$.
The first inequality follows from Jensen's inequality for the convex function $x
\mapsto e^{sx}$. The second inequality follows from the fact that maximum of
non-negative numbers is upper bounded by their sum. The last inequality follows from
the sub-Gaussian property for $Y_k$.

Thus, for any $s \in \R$,
$$
\exp\left(s \Exp\left[ \max_{k=1,2,\dots,K} |Y_k| \right] \right) \le 2K e^{s^2\sigma^2/2} \: .
$$
Taking logarithms, we get that for any $s \in \R$,
$$
s \Exp\left[ \max_{k=1,2,\dots,K} |Y_k| \right] \le \frac{s^2 \sigma^2}{2} + \ln(2K) \: .
$$
We restrict ourselves to $s > 0$ and divide both sides by $s$. We get that for any $s > 0$,
$$
\Exp\left[ \max_{k=1,2,\dots,K} |Y_k| \right] \le \frac{s \sigma^2}{2} + \frac{\ln(2K)}{s} \: .
$$
If $\sigma > 0$, we choose $s = \frac{2\sqrt{\ln(2K)}}{\sigma}$ and get that
$$
\Exp\left[ \max_{k=1,2,\dots,K} |Y_k| \right] \le \frac{s \sigma^2}{2} + \frac{\ln(2K)}{s} = \sigma \sqrt{2 \ln(2K)} \: .
$$
If $\sigma = 0$ then for any $s > 0$,
$$
\Exp\left[ \max_{k=1,2,\dots,K} |Y_k| \right] \le \frac{\ln(2K)}{s} \: .
$$
Taking $s \to \infty$, we get that $\Exp\left[ \max_{k=1,2,\dots,K} |Y_k| \right] \le 0$.
\end{proof}

\end{document}
